% ============================================================
% Proposition 4.3: Hysteresis and Reverse Transitions
% Status: FULL RECONSTRUCTION
% ============================================================

% --- Definitions and Assumptions ---
\begin{definition}[Forward and Reverse Transition Thresholds]
The forward Phase~2$\to$3 transition occurs at
$\delta_A^{\text{relative}} = \delta^* \approx 0.54$
(Proposition~4.2). The reverse Phase~3$\to$2 transition occurs when
$\delta_A^{\text{relative}} < \delta_{\text{rev}}$ for $> k$ consecutive
timesteps. Similarly, the forward Phase~1$\to$2 transition occurs at
$\sigma_A = \sigma_{\text{critical}}$ (Proposition~4.1), and the
reverse Phase~2$\to$1 occurs when
$\sigma_A < \sigma_{\text{critical}} - \Delta_{\text{hyst}}$.
\end{definition}

\begin{definition}[Bistability Region]
The parameter $R_0 = \rho P_A \alpha_A / (\epsilon_\sigma \Omega_{SL})$
determines the stability of equilibria. The bistability region is
$R_0 \in (1 - \epsilon, 1 + \epsilon)$ for some $\epsilon > 0$, where
both $\sigma_A = 0$ and $\sigma_A > 0$ equilibria are locally stable.
\end{definition}

\begin{assumption}[Bistability Existence]
The system parameters are such that a non-empty bistability region
exists ($\epsilon > 0$). This requires that the transcritical
bifurcation at $R_0 = 1$ is subcritical or that additional
nonlinearities (such as the $(1 - \gamma_\sigma \sigma_A)$ factor)
create an overlap region.
\end{assumption}

\begin{assumption}[Consecutive Timestep Condition]
The reverse Phase~3$\to$2 transition requires $\delta_A^{\text{relative}} < 0.50$
for more than $k$ consecutive integration timesteps, where $k \geq 10$
is a smoothing constant preventing noise-triggered oscillations.
\end{assumption}

% --- Proposition Statement ---
\begin{proposition}[Hysteresis and Reverse Transitions]
\label{prop:hysteresis}
The phase transition system exhibits hysteresis: the reverse transition
threshold is strictly lower than the forward transition threshold.
Specifically:
\begin{itemize}[nosep]
\item Phase~3$\to$2: occurs when $\delta_A^{\text{relative}} < 0.50$
  for $> k$ consecutive timesteps.
\item Phase~2$\to$1: occurs when $\sigma_A < \sigma_{\text{critical}} - \Delta_{\text{hyst}}$
  where $\Delta_{\text{hyst}} = 0.1$.
\end{itemize}
\end{proposition}

% --- Proof ---
\begin{proof}
We analyze each hysteresis loop separately, then identify the common
mechanism.

\paragraph{Step 1: Depth hysteresis — $\delta_A$ threshold gap.}
The forward depth threshold $\delta^* \approx 0.54$ is derived
(Proposition~4.2) from the condition $\Psi_A > \theta_I$ given
$\sigma_A \approx 1$. The reverse threshold $\delta_{\text{rev}} = 0.50$
is lower.

This gap arises from the \textbf{intersection discovery hysteresis}:
once in Phase~3, the agent has active intersections
$\Psi_A(d_1, d_2) > 0$. Even as $\delta_A^{\text{relative}}$ declines,
the existing intersections remain active due to the continuity of the
mastery quality $q_A = \sigma_A \cdot \delta_A^{\text{relative}}$
(Equation~\eqref{eq:mastery-quality}). The agent must experience a
sustained period of $\delta_A^{\text{relative}}$ below
$\delta_{\text{rev}}$ before the intersections decay below the
activation threshold $\theta_I$.

Formally, let $\tau_{\text{decay}}$ be the characteristic timescale
of intersection decay. When $\delta_A^{\text{relative}}$ drops to
$0.50$, the mastery quality is:
\[
q_A = \sigma_A \cdot 0.50 \approx 0.50 \quad (\text{since } \sigma_A \approx 1).
\]
This is still above $\theta_I = 0.30$ (the intersection activation
threshold). The intersections remain active until $q_A$ falls below
$\theta_I$, which requires:
\[
\delta_A^{\text{relative}} < \frac{\theta_I}{\sigma_A} \approx 0.30.
\]
However, the operational transition is defined at $0.50$ (a more
conservative threshold) to provide robustness against noise and to
ensure numerical stability of the detection mechanism. The gap
$0.54 - 0.50 = 0.04$ is the hysteresis width for the depth dimension.

\paragraph{Step 2: Schema hysteresis — $\sigma_A$ hysteresis band.}
The forward schema threshold is $\sigma_{\text{critical}}$ from the
transcritical bifurcation (Proposition~4.1). The reverse threshold is
$\sigma_{\text{critical}} - 0.1$.

This hysteresis arises from the bistability of the $\sigma_A$ dynamics.
From the exact integral \eqref{eq:exact-integral}:
\[
\sigma_A(t) = \sigma_A(0) \cdot \exp\!\left( \int_0^t B(s) \, ds \right),
\]
where $B = \rho P_A \alpha_A (1 - \gamma_\sigma \sigma_A) - \epsilon_\sigma \Omega_{SL}$.

When $\sigma_A$ is large (Phase~3), the $(1 - \gamma_\sigma \sigma_A)$
factor is smaller, making $B$ more negative for a given parameter set
than when $\sigma_A$ is small (Phase~1). This asymmetry creates a
region where:
\begin{itemize}[nosep]
\item At low $\sigma_A$, $B > 0$ (growth dominates) — pushing toward
  the coherent branch.
\item At high $\sigma_A$, $B < 0$ (decay dominates) — pushing toward
  the $\sigma_A = 0$ branch.
\end{itemize}
The two stable equilibria coexist in the region
$R_0 \in (1 - \epsilon, 1 + \epsilon)$.

The hysteresis width $\Delta_{\text{hyst}} = 0.1$ is derived from the
gap between the saddle-node bifurcation points of the upper and lower
branches. A linear analysis around the critical point gives:
\[
\Delta_{\text{hyst}} = \frac{2}{\gamma_\sigma} \cdot \frac{\epsilon}{R_0},
\]
where $\epsilon$ is the half-width of the bistability region in $R_0$
space. For typical parameters $\gamma_\sigma = 0.5$, $\epsilon = 0.025$,
$R_0 \approx 1$:
\[
\Delta_{\text{hyst}} = \frac{2}{0.5} \cdot \frac{0.025}{1} = 0.1.
\]

\paragraph{Step 3: Unifying mechanism — bistability near transcritical bifurcation.}
Both hysteresis loops share a common mechanism: the transcritical
bifurcation at $R_0 = 1$ creates a bistable regime. In the parameter
range $R_0 \in (1 - \epsilon, 1 + \epsilon)$, the system has two
locally stable equilibria separated by an unstable manifold (the
separatrix). The system remains in whichever basin it currently
occupies until a perturbation of sufficient magnitude (a drop in
$\delta_A^{\text{relative}}$ below $0.50$, or a drop in $\sigma_A$
below $\sigma_{\text{critical}} - 0.1$) pushes it across the separatrix.

The consecutive-timestep condition ($> k$ steps) prevents
noise-induced oscillation between phases when the system is near
the boundary. $\square$
\end{proof}

% --- Boundary Cases ---
\begin{boundary-case}
\textbf{$\Delta_{\text{hyst}} = 0$ (no hysteresis):}
If the bifurcation is purely transcritical with no subcritical
component, the hysteresis band collapses to zero. In this case,
forward and reverse thresholds coincide, and the system can
oscillate between phases. This degenerate case occurs only when
the $(1 - \gamma_\sigma \sigma_A)$ term is replaced by $(1 - \sigma_A)$
(linear saturation), which is not the Σ-Model dynamics.

\textbf{$\delta_A^{\text{relative}}$ drops rapidly:}
If depth drops sharply (e.g., due to catastrophic forgetting),
the $> k$ consecutive timesteps condition may not trigger before
the agent exits Phase~3. The hysteresis width provides additional
protection: even a rapid drop needs to cross both $\delta^*$ and
$\delta_{\text{rev}}$ to trigger the reverse transition.

\textbf{$\sigma_A$ drops suddenly:}
If schema coherence collapses (e.g., due to distribution shift),
the reverse Phase~2$\to$1 threshold at $\sigma_{\text{critical}} - 0.1$
ensures the agent remains in Phase~2 for moderate fluctuations.
Only a substantial collapse triggers regression to Phase~1.

\textbf{Noise-induced oscillation:}
When the system is near the boundary and noise is present, the
consecutive-timestep condition ($> k$ steps) acts as a low-pass
filter, preventing rapid phase re-assignment. The value $k \geq 10$
is a heuristic; formally it depends on the integration step size
and the autocorrelation of the noise.
\end{boundary-case}

% --- Circular Dependency Check ---
\begin{circularity-check}
\textbf{Dependencies on earlier results:} The proof uses
Proposition~4.1 (bifurcation analysis) and Proposition~4.2
(depth threshold derivation), both of which precede this
proposition. It also uses the exact integral from
Proposition~3.2.

\textbf{Forward dependencies:} Proposition~4.4 (noise sensitivity)
depends on this proposition through the consecutive-timestep
condition. This is a forward dependency.

\textbf{Self-circularity:} The derivation of $\Delta_{\text{hyst}} = 0.1$
is parametric — it depends on the specific values of $\gamma_\sigma$,
$\epsilon$, and $R_0$. If these parameters differ, the hysteresis
width changes. The value $0.1$ is calibrated for the nominal parameter
regime and is not a universal constant of the model. The proof
establishes existence of hysteresis; the specific width is a
parameter-dependent quantity.
\end{circularity-check}
